\chapter{Introduction}
\label{sec:intro}

This document answers two questions we posed ourselves while planning a
shared tuning system for the projects in our workspace: whether Ray Tune
is the right architecture to build on, given that RLlib already runs our
reinforcement-learning work, and which algorithms such a system should
implement to serve very different customers, from reinforcement-learning
agents through Hamiltonian neural networks and neural-transport samplers to
multi-objective models in finance. Both questions turn out to have
reasonably confident answers once the literature is laid out, and laying
it out carefully is most of the document: our local paper collection
thins out around 2023, and several things a 2023 reader would recommend
are no longer what a 2026 reader should build.

The short answers, defended over the following chapters; the method
names that follow are only names for now, each taught from zero in its
chapter, so read this section as a map to return to rather than an
argument to absorb. Ray Tune: yes,
as the execution layer, and deliberately not as the source of search
algorithms. Tune's distributed trial execution, checkpoint machinery, and
population-based training support remain the strongest in open source and
compose naturally with RLlib, but its own algorithm layer has been in
maintenance for roughly two years while the method frontier moved into
other projects (Chapters~\ref{sec:libraries}
and~\ref{sec:architecture}).

The algorithms: a spine of asynchronous successive halving plus a
well-configured Bayesian optimizer covers the routine work at low cost;
the differentiators worth building are noise-aware multi-objective
acquisition with a multi-objective early-stopping scheduler, expert
priors, cost-aware acquisition, and a population line for reinforcement
learning. The method names that follow are only names for now, each
taught from zero in its chapter: qLogNEHVI with MO-ASHA
(Chapters~\ref{sec:moo}), priors in the PriorBand style
(Chapter~\ref{sec:transfer}), and the population line as PB2 extended to
mixed spaces and then BG-PBT, the method we were specifically asked to
assess (Chapter~\ref{sec:population}). Freeze-thaw scheduling with a
pre-trained learning-curve surrogate is the one frontier bet we recommend
(Chapter~\ref{sec:multifidelity}), and Chapter~\ref{sec:roadmap} puts
the whole list in build order with the evidence attached to each item.

The document is written to teach as well as to conclude, because the
people extending this system next year will need the reasoning, not just
the verdicts, and a verdict whose derivation you cannot reconstruct is
one you cannot safely modify. Three habits follow, and knowing them up
front makes the document easier to read.

First, one running example
carries the whole text: tuning PPO on a Brax locomotion task, in the
fifteen-dimensional search space introduced in
Section~\ref{sec:running}, with a Hamiltonian network as the second
thread wherever several objectives compete. When a formula arrives, its
symbols are read against that example before anything is done with
them.

Second, every load-bearing equation is derived or unpacked in
place rather than quoted: the Gaussian-process posterior, expected
improvement and the floating-point failure that motivates its modern
replacement, the tree-structured Parzen estimator's ratio rule, the
budget arithmetic of successive halving, the geometry that breaks
weighted-sum scalarization.

Third, the figures are schematic drawings
of mechanisms, not plots of measured data, and their captions say what
each one is supposed to show; when a number matters, it is computed in
the text where the reader can check it.

The document also carries a decision layer, kept visually separate from
the teaching so that neither gets in the other's way. Each method
chapter ends in a boxed \emph{Decisions from this chapter} marked by a
thick left rule: plain-language verdicts, include or exclude, with the
cost to build, the evidence behind the call, and the comparison against
the alternatives, written to be judged without opening any cited paper.

The box is earned rather than asserted: each method chapter closes by
first teaching the experiments its verdict rests on, their designs,
arenas, and sponsors, so the evidence can be judged rather than taken
on citation, and then arguing the verdict from several independent
angles, ending with the criterion that would overturn it in our own
validation suite.

Table~\ref{tab:decisions} in Chapter~\ref{sec:roadmap} collects every
verdict on one page. A glossary at the end defines every term of art in
plain language. And because several of our readers come from economics
and statistics rather than machine learning, the text flags the bridges
where they exist: the Gaussian process is kriging, expected improvement
is an option value, the Pareto front (the set of configurations
undominated in objective space) is the welfare-theoretic one.

Chapter~\ref{sec:problem} pins down the problem, fixes the running
example, and maps the four ways our workloads bend the textbook
formulation. Chapters~\ref{sec:modelfree} through~\ref{sec:transfer}
build up the method families in increasing order of structure: searching
without a model, modeling the objective, exploiting partial evaluations,
tuning while training, handling several objectives, and importing
knowledge from priors, past studies, and scale. Each chapter opens with
the failure of the previous chapter's best answer, which is not a
storytelling flourish: it is the actual dependency structure of the
field, and it is why the chapters are in this order.

Chapter~\ref{sec:workloads} turns back to our four workload families and
says what to do for each. Chapter~\ref{sec:libraries} audits the
open-source field library by library, with maintenance status checked
against primary sources on August 22, 2026, and re-verified on
August 23. The closing part turns the verdicts into a concrete product
proposal: Chapter~\ref{sec:product} states what we ship and walks a
first study through it, Chapter~\ref{sec:architecture} gives the
architectural design, Chapter~\ref{sec:contracts} pins the interfaces
and data schemas, Chapter~\ref{sec:validation} specifies the test
suite, and Chapter~\ref{sec:roadmap} the build plan and the bill.

A reader who wants
only decisions can read Chapter~\ref{sec:problem}, the decision boxes
at the ends of Chapters~\ref{sec:modelfree}
through~\ref{sec:workloads}, and then Chapters~\ref{sec:libraries},
\ref{sec:architecture}, and~\ref{sec:roadmap} with
Table~\ref{tab:decisions}; a reader implementing a searcher should
read the method chapters in order, since each one leans on the
machinery of its predecessors.

A note on sourcing. Claims about specific methods are tied to the papers
that introduced or tested them, and the load-bearing ones (BG-PBT, the
RL tuning study of Eimer, Lindauer, and Raileanu, and the current
defaults of the major libraries) were checked against the primary texts
and pages rather than summaries; a companion file records, for every
source, what was inspected and at what depth. The derivations added for
teaching (Gaussian conditioning, the expected-improvement closed form
and its underflow arithmetic, the successive-halving budget count, the
hypervolume and scalarization geometry) are elementary and stand on
their own; they explain the cited results rather than extend them.
Where a comparison rests on one paper's experiments we say so, and
where the honest answer is that two methods are indistinguishable on
current evidence, we say that too, in preference to manufacturing a
ranking.
